\documentclass[french]{article}

\usepackage{babel}
\usepackage{comment}
\usepackage{psfig}
\newenvironment{solution}{\nopagebreak\par\vspace*{2mm}\underline{Solution~:}\em\par}{\par\vspace*{2mm}}
\usepackage{times}
\usepackage[latin1]{inputenc}


\pagestyle{headings}
\markright{28/6/1995 \hfill Conservatoire National des Arts et
M\'etiers -- Paris \hfill\ } 

\baselineskip 6pt
\parskip 6pt
\onecolumn
\textwidth 16cm
\textheight 24cm
\hoffset=-2cm
\voffset=-2.5cm
\setlength{\parindent}{0.0cm}
\unitlength 1cm
%\input{psfig}
\newtheorem{defin}{D\'efinition}[section]
\newtheorem{exampl}{Exemple}[section]
\def\picwidth{16cm}

%\excludecomment{solution}

\begin{document}
\begin{center}
\Large\bf Processus d'Informatisation et Bases de Donn\'ees A \\ 
    Examen Bases de Donn\'ees\\
\large\bf Samedi 16 septembre 2000
\end{center}

%%%%%%%%%%%%%%%%%%%%%%%%%%%%%%%%%%%%%%%%%%%%%%%%%%%%%%%%%%%%

\vspace*{0.5cm}

Une startup d�cide de cr�er un site web proposant tous les restaurants
parisiens. Le site est construit sur une base de donn�es
contenant les restaurants, les plats propos�s et les menus.
Voici les tables.  Comme d'habitude, les cl�s �trang�res 
ne sont pas repr�sent�es\,:
      \begin{itemize}
        \item {\tt Restaurant (\underline{nomRestau}, arrondissement)}
         \item {\tt Plat (\underline{nomRestau, nomPlat}, prix, typePlat)}
         \item {\tt Menu (\underline{nomRestau, noMenu}, nom, prix)}
         \item {\tt D�taiMenu (\underline{nomRestau, noMenu, nomPlat})} 
       \end{itemize}

L'attribut {\tt arrondissement} contient le num�ro
de l'arrondissement du restaurant.
La table {\tt Plat} donne l'ensemble des plats propos�s par un restaurant.
Le type du plat est soit {\tt Entr�e}, soit {\tt Plat principal},
soit {\tt Dessert}.
La table {\tt D�tailMenu} donne les plats composant un menu\,:
un ensemble d'entr�es, de plats principaux et de dessert.

\smallskip
{\bf Requ\^etes ({\bf 10 pts})}

Exprimez en SQL les requ\^etes suivantes:

\begin{enumerate}
  \item  ({\bf 1 pt}) Nom des restaurants du douzi�me arrondissement.

  \item  ({\bf 1 pt}) Nom des restaurants dont le nom
      contient le mot <<\,fourchette\,>>.

  \item  ({\bf 2 pts})  Quels 
               restaurants (nom et arrondissement)
                 proposent des menus � 150 frs\,?
  
  \item  ({\bf 2 pts}) Donnez la liste des entr�es du menu num�ro 2 du restaurant
    <<\,Boboche\,>>.

  \item  ({\bf 2 pts}) Donnez les menus (nomRestau, noMenu)
               ayant un plat plus cher que le prix du menu lui-m�me.

  \item  ({\bf 2 pts}) Quels restaurants (nomRestau) ont un menu qui ne propose
    pas de dessert\,?

  \end{enumerate}

\smallskip
{\bf Sch�ma (10 pts)}

\begin{enumerate}
   \item ({\bf 3 pts}) Donnez les commandes {\tt CREATE TABLE} d\'ecrivant
         le sch\'ema des relations {\tt Plat}, {\tt Menu} et
            {\tt DetailMenu}. Indiquez 
         les cl\'es \'etrang\`eres et les cl\'es primaires.

   \item ({\bf 1 pts}) Donnez la contrainte SQL2 qui force
      le type d'un plat aux trois valeurs possibles.
 
   \item ({\bf 2 pts})
       On veut permettre aux internautes de noter les restaurants.
       Pour cela on veut conna�tre les internautes (nom, pr�nom, email)
       et savoir que tel internaute a donn� telle note � tel restaurant.

       Donnez le sch�ma E/A correspondant � l'ajout de ces informations.

   \item  ({\bf 2 pts}) Donnez le sch�ma des nouvelles tables � cr�er.

    \item ({\bf 2 pts})
           Quel est l'ordre SQL donnant la moyenne des notes obtenues
        par <<\,Boboche\,>>\,? 
\end{enumerate}
    
\end{document}
